\section{Error-Correcting Codes}
\definecolor{accent}{HTML}{1F4E79}
\definecolor{softaccent}{HTML}{EAF2F8}
\definecolor{softgray}{HTML}{F6F7F9}
\definecolor{successdraw}{HTML}{27AE60}
\definecolor{successfill}{HTML}{EAFAF1}
\definecolor{errordraw}{HTML}{C0392B}
\definecolor{errorfill}{HTML}{FDEDEC}

The fundamental problem of coding theory is to transmit information reliably across a noisy channel. 
A sender begins with a message, usually represented as a vector
\[
\mathbf{x}=(x_1,\dots,x_k)\in \F_2^k,
\]
and applies an \emph{encoder}
\[
\Encode:\F_2^k \longrightarrow \F_2^n,
\qquad
\mathbf{x}\longmapsto \mathbf{c},
\]
where \(n\ge k\). 
The vector \(\mathbf{c}\) is called the \emph{codeword}. 
%The passage from length \(k\) to length \(n\) introduces redundancy, and this redundancy is precisely what makes error detection and error correction possible.

During transmission, the codeword passes through a noisy channel. 
The effect of the channel is modeled by adding an error vector $
\mathbf{e}\in \F_2^n$
to the transmitted codeword, so that the receiver obtains
\[
\mathbf{y}=\mathbf{c}\oplus\mathbf{e}\in\F_2^n.
\]
The vector \(\mathbf{y}\) is called the \emph{received vector}. 
The purpose of the decoder is to recover the original message, or equivalently the transmitted codeword, from the corrupted vector \(\mathbf{y}\).

\begin{center}
	\begin{tikzpicture}[
	>=Stealth,
	block/.style={
		draw=accent, thick, fill=softaccent, rounded corners=2mm,
		minimum width=1.2cm, minimum height=1cm, align=center,
		font=\sffamily, drop shadow={opacity=0.15, shadow xshift=1mm, shadow yshift=-1mm}
	},
	adder/.style={
		circle, draw=accent, thick, fill=softaccent,
		minimum size=6mm, inner sep=0pt, align=center, font=\Large,
		drop shadow={opacity=0.15, shadow xshift=1mm, shadow yshift=-1mm}
	},
	arrow/.style={->, thick, accent},
	mathlabel/.style={font=\normalsize, text=black},
	sublabel/.style={font=\small\sffamily, text=magenta},
	scale=.85
	]
	
	% Invisible nodes to act as the start and end points of the signals
	\node (start) at (0,0) {};
	\node (end)   at (19,0) {};
	
	% Main operational blocks
	\node[block, align=center] (enc) at (4.5,0) {Encoder\\ $\F_2^k\to\F_2^n$};
	\node[adder] (add) at (9.5,0) {$+$};
	\node[block, align=center] (dec) at (14.5,0) {Decoder\\ $\F_2^n\to\F_2^k$};
	
	% Noise entry point
	\node (noise) at (9.5, 2) {\small \texttt{Noise} \textbf{e}};
	
	% --- Signal Paths & Labels ---
	
	% Source to Encoder
	\draw[arrow] (start) -- 
	node[midway, mathlabel, draw=accent, fill=white, align=center] {\texttt{message}\\ [-2.5mm] $\textbf{x}$} 
	node[below=5mm, sublabel] {$k$-bit} 
	(enc);
	
	% Encoder to Channel (Adder)
	% We use a standard label below the arrow to replace the underbrace in the equation
	\draw[arrow] (enc) -- 
	node[midway, mathlabel, draw=accent, fill=white, align=center] {\texttt{codeword}\\ [-2.5mm] $\textbf{c}$} 
	node[below=5mm, sublabel] {$n$-bit ($n \geq k$)} 
	(add);
	
	% Noise entering the Channel
	\draw[arrow] (noise) -- (add);
	
	% Channel to Decoder
	\draw[arrow] (add) -- 
	node[midway, mathlabel, draw=accent, fill=white, align=center] {\texttt{received}\\ [-2.5mm]  \texttt{vector} $\textbf{y}$} 
	node[below=5mm, sublabel] {$n$-bit} 
	(dec);
	
	% Decoder to Destination
	\draw[arrow] (dec) -- 
	node[midway, mathlabel, draw=accent, fill=white, align=center] {\texttt{message}\\ [-2.5mm] $\widehat{\textbf{x}}$} 
	node[below=5mm, sublabel] {$k$-bit} 
	(end);
	
	\draw[dashed, red] (8.6,-1.2) rectangle (10.4,2.5);
	\node[red] at (9.5,3) {channel};
\end{tikzpicture}
\end{center}
%Thus the overall communication process is modeled by
%\[
%\mathbf{x}
%\;\xmapsto{\ \Encode\ }\;
%\mathbf{c}
%\;\xmapsto{\ \text{channel}\ }\;
%\mathbf{y}=\mathbf{c}+\mathbf{e}
%\;\xmapsto{\ \Decode\ }\;
%\widehat{\mathbf{x}}.
%\]
%A good error-correcting code is one for which the decoder can still recover the original message whenever the error vector is not too large, in a sense measured by its Hamming weight. 
%In particular, if the code is designed to correct up to \(t\) errors, then one expects
%\[
%\Decode(\mathbf{y})=\mathbf{x}
%\qquad\text{whenever}\qquad
%\wt(\mathbf{e})\le t.
%\]
A generic communication model has the form
\[
\vect{x}=(x_1,\dots,x_k)
\xrightarrow{\text{encoder}}
\underbrace{\vect{c}=(c_1,\dots,c_n)}_{\text{codeword}\; (n\geq k)}
\xrightarrow{\text{noise }\vect e}
\vect{y}=\vect{c}+\vect{e}
\xrightarrow{\text{decoder}}
\widehat{\vect{x}}=(\widehat x_1,\dots,\widehat x_k).
\]

\medskip

\begin{center}
\begin{tikzpicture}[scale=1,>=Stealth,
	block/.style={
		draw=accent,
		thick,
		fill=softaccent,
		rounded corners=2mm,
		minimum width=3.2cm,
		minimum height=1cm,
		align=center,
		font=\sffamily,
		drop shadow={opacity=0.15, shadow xshift=1mm, shadow yshift=-1mm}
	},
	decision/.style={
		diamond,
		aspect=1.5, % Stretches the diamond horizontally to fit text
		draw=accent,
		thick,
		fill=softgray,
		align=center,
		font=\sffamily\small,
		drop shadow={opacity=0.15, shadow xshift=1mm, shadow yshift=-1mm}
	},
	successblock/.style={
		block,
		draw=successdraw,
		fill=successfill,
		text=black!85
	},
	errorblock/.style={
		block,
		draw=errordraw,
		fill=errorfill,
		text=black!85
	},
	arrow/.style={
		->,
		thick,
		accent,
		rounded corners=2mm 
	},
	dashedarrow/.style={
		->,
		thick,
		accent,
		dashed, % Distinguishes evaluation logic from data flow
		rounded corners=2mm 
	}]
	
	% --- Main Pipeline Nodes ---
	\node[block] (tx) at (0,0) {Data $\vect{x}$};
	\node[block] (enc) at (6,0) {Encoder};
	\node[block] (chan) at (10,-1) {Channel + noise};
	\node[block] (dec) at (6,-2) {Decoder};
	\node[block] (rx) at (0,-2) {Data $\hat{\vect{x}}$};
	
	% --- Evaluation Nodes ---
%	\node[decision] (comp) at (0,-6) {Is $\hat{\vect{x}} = \vect{x}?$};
%	\node[successblock] (succ) at (6,-5) {\textbf{Successful Decoding}\\$\hat{\vect{x}} = \vect{x}$};
%	\node[errorblock] (err) at (6,-7) {\textbf{Decoding Error}\\$\hat{\vect{x}} \neq \vect{x}$};
	
	% --- Main Pipeline Arrows ---
	\draw[arrow] (tx) -- node[above, font=\small] {$\vect{x}\in\F_q^k$} (enc);
	\draw[arrow] (enc) -| node[above right, font=\small, pos=0.25] {$\vect{c}\in\F_q^n$} (chan);
	\draw[arrow] (chan) |- node[below right, font=\small, pos=0.75] {$\vect{y}=\vect{c}+\vect{e}\in\F_q^n$} (dec);
	\draw[arrow] (dec) -- node[above, font=\small] {$\widehat{\vect{x}}\in\F_q^k$} (rx);
	
	% --- Evaluation Arrows ---
%	% Route from estimated data down to comparison
%	\draw[dashedarrow] (rx.south) -- node[right, font=\small] {Estimate} (comp.north);
	
%	% Route from original data around the left side down to comparison
%	\draw[dashedarrow] (tx.west) -- node[above, font=\small] {Original} ++(-1.5,0) |- (comp.west);
	
%	% Branching arrows from comparison to success/error
%	\draw[arrow] (comp.east) -- ++(1.5,0) |- (succ.west);
%	\draw[arrow] (comp.east) -- ++(1.5,0) |- (err.west);
\end{tikzpicture}
\end{center}
\noindent


\newpage

\begin{example}[Repetition code]
	Let \(\Sigma\) be an alphabet. An \emph{encoding} is a map
	\[
	\fullfunction{\mathsf{Encode}}{\Sigma^k}{\Sigma^n}{(x_1,\dots,x_k)}{(c_1,\dots,c_n)}
	\] with \(n\ge k\). %The output vector \(\vect c=(c_1,\dots,c_n)\) is the transmitted codeword.
	\begin{center}
		\begin{tikzpicture}[
			>=Stealth,
			% Style for single input symbols
			symbol/.style={
				circle, draw=accent, thick, fill=softgray,
				minimum size=9mm, inner sep=0pt, align=center, font=\normalsize,
				drop shadow={opacity=0.15, shadow xshift=1mm, shadow yshift=-1mm}
			},
			% Style for the output codeword arrays
			block/.style={
				rectangle, draw=accent, thick, fill=softaccent, rounded corners=2mm,
				minimum width=3.8cm, minimum height=1.1cm, align=center, font=\normalsize,
				drop shadow={opacity=0.15, shadow xshift=1mm, shadow yshift=-1mm}
			},
			% Style for the mapping arrows
			maparrow/.style={
				|->, thick, accent, shorten >=1.5mm, shorten <=1.5mm
			},
			labeltext/.style={
				font=\small\sffamily, text=black!85
			}
			]
			% --- General Formula Row ---
			\node[symbol] (x1) at (0, 0) {$x_1$};
			\node[block] (outX) at (5.5, 0) {$(x_1, \, x_1, \, \dots, \, x_1)$};
			
			\draw[maparrow] (x1) -- node[above, labeltext] {$\mathsf{Encode}$} (outX.west);
			
			% Brace under the general codeword to show length
			\draw[thick, accent, decorate, decoration={brace, amplitude=5pt, mirror}]
			([xshift=4mm, yshift=-2.5mm]outX.south west) -- 
			([xshift=-4mm, yshift=-2.5mm]outX.south east)
			node[midway, below=6pt, labeltext] {$n$ copies ($\in \Sigma^n$)};
		\end{tikzpicture}
	\end{center}
	%If \(\Sigma=\{0,1\}\) then
	%\[
	%(0)\mapsto (0,0,\dots,0),
	%\qquad
	%(1)\mapsto (1,1,\dots,1).
	%\]
	When \(k=1\), a very simple encoding repeats the same symbol
	%\[
	%(x_1)\longmapsto (x_1,x_1,\dots,x_1)\in \Sigma^n.
	%\]
	\begin{center}
		\begin{tikzpicture}[
			>=Stealth,
			% Style for single input symbols
			symbol/.style={
				circle, draw=accent, thick, fill=softgray,
				minimum size=9mm, inner sep=0pt, align=center, font=\normalsize,
				drop shadow={opacity=0.15, shadow xshift=1mm, shadow yshift=-1mm}
			},
			% Style for the output codeword arrays
			block/.style={
				rectangle, draw=accent, thick, fill=softaccent, rounded corners=2mm,
				minimum width=3.8cm, minimum height=1.1cm, align=center, font=\normalsize,
				drop shadow={opacity=0.15, shadow xshift=1mm, shadow yshift=-1mm}
			},
			% Style for the mapping arrows
			maparrow/.style={
				|->, thick, accent, shorten >=1.5mm, shorten <=1.5mm
			},
			labeltext/.style={
				font=\small\sffamily, text=black!85
			},
			header/.style={
				font=\bfseries\sffamily, text=accent
			}
			]
			% --- Binary Case: Input 0 ---
			% Using white fill to contrast with the abstract case
			\node[symbol, fill=white] (b0) at (0, -2.8) {$0$};
			\node[block, fill=white] (out0) at (5.5, -2.8) {$(0, \, 0, \, \dots, \, 0)$};
			
			\draw[maparrow] (b0) -- (out0.west);
			
			% --- Binary Case: Input 1 ---
			\node[symbol, fill=white] (b1) at (0, -4.3) {$1$};
			\node[block, fill=white] (out1) at (5.5, -4.3) {$(1, \, 1, \, \dots, \, 1)$};
			
			\draw[maparrow] (b1) -- (out1.west);
		\end{tikzpicture}
	\end{center}
	This is called the \emph{repetition code}. If one symbol is changed during transmission, we can still recover the original message by \emph{majority vote}. For instance,
	\[
	(0,0,0)\rightsquigarrow (0,1,0),
	\qquad
	(1,1,1)\rightsquigarrow (1,0,1).
	\]
	Since \(0\) appears most often in \((0,1,0)\), we decode it as \(0\), and since
	\(1\) appears most often in \((1,0,1)\), we decode it as \(1\).
	%More generally, if \(n\) is odd, then the repetition code can correct up to
	%\[
	%\left\lfloor \frac{n-1}{2}\right\rfloor
	%\]
	%errors, because the original symbol is still the majority symbol.
	
	\medskip
	
	\begin{center}
		\definecolor{success}{HTML}{27AE60}
		\begin{tikzpicture}[
			>=Stealth,
			block/.style={
				rectangle, draw=accent, thick, fill=softaccent, rounded corners=2mm,
				minimum width=2.2cm, minimum height=1cm, align=center, font=\large,
				drop shadow={opacity=0.15, shadow xshift=1mm, shadow yshift=-1mm}
			},
			decision/.style={
				rectangle, draw=accent, thick, fill=softgray, rounded corners=2mm,
				minimum width=3.2cm, minimum height=1.1cm, align=center, font=\sffamily\small,
				drop shadow={opacity=0.15, shadow xshift=1mm, shadow yshift=-1mm}
			},
			result/.style={
				circle, draw=success, thick, fill=success!10, 
				minimum size=10mm, inner sep=0pt, align=center, font=\Large\bfseries, text=success,
				drop shadow={opacity=0.15, shadow xshift=1mm, shadow yshift=-1mm}
			},
			arrow/.style={
				->, thick, accent
			},
			noisearrow/.style={
				->, thick, errordraw, 
				% The snake decoration mimics the \rightsquigarrow command
				decorate, decoration={snake, amplitude=0.4mm, segment length=3mm, post length=1.5mm}
			},
			labeltext/.style={font=\small\sffamily, text=black!85},
			header/.style={font=\bfseries\sffamily, text=accent}
			]
			
			% --- Column Headers ---
			\node[header] at (0, 1.5) {Original Data};
			\node[header] at (4, 1.5) {Received};
			\node[header] at (8.5, 1.5) {Decoder};
			\node[header] at (12.8, 1.5) {Recovered Data};
			
			\node[block] (tx1) at (0, 0) {$0$};
			
			% The corrupted bit is highlighted in bold red
			\node[block] (rx1) at (4, 0) {$(0, \textcolor{errordraw}{\mathbf{1}}, 0)$};
			
			\node[decision] (vote1) at (8.5, 0) {Two 0s, One 1\\[1mm] (\textbf{0 is the majority})};
			\node[result] (out1) at (12.8, 0) {0};
			
			\draw[noisearrow] (tx1) -- node[above, font=\small\sffamily, text=errordraw] {error} (rx1);
			\draw[arrow] (rx1) -- node[above, labeltext] {} (vote1);
			\draw[arrow] (vote1) -- (out1);
			
			% --- Horizontal Divider ---
			%	\draw[dashed, draw=gray!40, thick] (-1.5, -1.6) -- (14, -1.6);
			
			%	% ==========================================
			% EXAMPLE 2: Sending 1s
			% ==========================================
			\node[block] (tx2) at (0, -3.2) {$1$};
			
			% The corrupted bit is highlighted in bold red
			\node[block] (rx2) at (4, -3.2) {$(1, \textcolor{errordraw}{\mathbf{0}}, 1)$};
			
			\node[decision] (vote2) at (8.5, -3.2) {Two 1s, One 0\\[1mm] (\textbf{1 is the majority})};
			\node[result] (out2) at (12.8, -3.2) {1};
			
			\draw[noisearrow] (tx2) -- node[above, font=\small\sffamily, text=errordraw] {error} (rx2);
			\draw[arrow] (rx2) -- node[above, labeltext] {} (vote2);
			\draw[arrow] (vote2) -- (out2);
		\end{tikzpicture}
	\end{center}
\end{example}

\newpage
\subsection{Linear Codes}
\defbox[Linear Encoding]{\begin{definition}
		A \textbf{linear encoding} over \(\F_q\) is a linear map
		\[
		\F_q^k \longrightarrow \F_q^n,
		\qquad
		\vect{x}\longmapsto \vect{c}.
		\]
\end{definition}}
%\begin{remark}
%Since the map is linear,
%\begin{itemize}
%	\item the sum of two codewords is again a codeword;
%	\item any scalar multiple of a codeword is again a codeword.
%\end{itemize}
%\end{remark}

\medskip

\defbox[Linear Code]{\begin{definition}
		%A linear subspace \(\code\subseteq \F_q^n\) is called a \textbf{linear code}.
		A \textbf{code of length \(n\in \mathbb{Z}_{>0}\) over \(\F_q\)} is a subset \[
		\code\subseteq \F_q^n.
		\]
		The elements of \(\code\) are called \textbf{codewords}. If \(\code\leq\F_q^n\), then \(\code\) is called a \textbf{linear code}.
\end{definition}}
\begin{center}
	\begin{tikzpicture}[
		>=Stealth,
		ambient/.style={
			draw=accent, thick, fill=softgray, rounded corners=3mm, 
			minimum width=5cm, minimum height=5cm,
			drop shadow={opacity=0.15, shadow xshift=1mm, shadow yshift=-1mm}
		},
		labeltext/.style={text=black!85},
		header/.style={font=\bfseries\sffamily\large, text=accent},
		codeword/.style={circle, fill=accent, inner sep=1.5pt}
		]
		
		% ==========================================
		% PANEL 1: General Code (Arbitrary Subset)
		% ==========================================
		
		% Ambient Space
		\node[ambient] (space1) at (0,0) {};
		\node[anchor=north west, labeltext] at (space1.north west) {\; Vector Space $\F_q^n$};
		
		% Arbitrary Subset (Blob)
		\draw[draw=accent, thick, fill=softaccent, rotate around={20:(0,0)}] 
		(0.2, -0.2) ellipse (1.6cm and 1cm) node[header, below=40pt] {Subset $\code\subseteq\F_q^n$};
		
		% Codewords inside the subset
		\node[codeword, label={[labeltext, font=\small]below:$\vect{c}_1$}] at (-0.5, -0.3) {};
		\node[codeword, label={[labeltext, font=\small]right:$\vect{c}_2$}] at (0.8, -0.6) {};
		\node[codeword, label={[labeltext, font=\small]above:$\vect{c}_3$}] at (0.6, 0.4) {};
		
		% The zero vector is explicitly OUTSIDE the code to show it's not required
		\node[circle, draw=accent, thick, inner sep=1.5pt, fill=white, label={[labeltext, font=\small]above:$\vect{0}$}] at (-1.5, 1) {};
		
		% ==========================================
		% PANEL 2: Linear Code (Vector Subspace)
		% ==========================================
		
		% Ambient Space
		\node[ambient] (space2) at (7,0) {};
		\node[anchor=north west, labeltext] at (space2.north west) {\; Vector Space $\F_q^n$};
		
		% Subspace (A plane passing through the origin)
		\draw[draw=accent, thick, fill=softaccent] 
		(4.75, -.9) -- (6.25, 1.5) -- (9.0, 1.2) -- (7.5, -1.2) -- cycle node[header, below right=20pt] {Subspace $\code\leq\F_q^n$};
		
		% The origin (MUST be in a linear code)
		\node[codeword, label={[labeltext, font=\small]below right:$\vect{0}$}] (zero) at (7,0) {};
		
		% Codewords demonstrating linear combinations
		\node[codeword, label={[labeltext, font=\small]left:$\vect{c}_1$}] (c1) at (6.2, 0.8) {};
		\node[codeword, label={[labeltext, font=\small]right:$\vect{c}_2$}] (c2) at (8.1, 0.3) {};
		
		% Vector addition (c1 + c2)
		\node[codeword, label={[labeltext, font=\small]above:$\vect{c}_1+\vect{c}_2$}] (c3) at (7.3, 1.1) {};
		
		% Dashed lines showing the parallelogram of vector addition
		\draw[dashed, accent, thin] (c1) -- (c3);
		\draw[dashed, accent, thin] (c2) -- (c3);
		\draw[->, accent, thin, shorten >=2pt] (zero) -- (c1);
		\draw[->, accent, thin, shorten >=2pt] (zero) -- (c2);
	\end{tikzpicture}
\end{center}

\begin{remark}
	If $\dim_{\F_q}(\code)=k$, then \(\code\) is called a \textbf{\([n,k]_q\)-code over \(\F_q\)}.
\end{remark}
%If \(\code\) is a linear code, then
%\begin{itemize}
%	\item the zero codeword \(\vect 0\) belongs to \(\code\);
%	\item whenever \(\vect c\in \code\), also \(-\vect c\in \code\);
%	\item every codeword is a linear combination of basis vectors of \(\code\).
%\end{itemize}

\probox[]{\begin{proposition}
If \(\code\) is an $\nkcode{n}{k}{q}$-code, then
\[
|\code|=q^k.
\]
\end{proposition}}
\begin{proof}
	Choose a basis \(\vect b_1,\dots,\vect b_k\) of \(\code\). Every codeword can be written uniquely as
	\[
	x_1\vect b_1+\cdots+x_k\vect b_k,
	\qquad x_i\in \F_q.
	\]
	Each \(x_i\) has \(q\) choices, and there are \(k\) independent coefficients. Therefore the number of codewords is \(q^k\).
\end{proof}

\newpage
\begin{flushleft}\color{blue}
\begin{note}[Linear Code as Short Exact Sequence]
A linear \([n,k]_q\) code is most naturally described by a short exact sequence \[
0 \longrightarrow \F_q^k \xrightarrow{\;\mathsf{Encode}\;} \F_q^n \xrightarrow{\;H\;} \F_q^{\,n-k} \longrightarrow 0,
%\qquad
%C=\operatorname{im}(\mathsf{Enc})=\ker(H).
\]
A message \(\vect x\in \F_q^k\) is encoded as
\[
\vect c=\mathsf{Enc}(\vect x)\in C.
\]
After transmission through the channel, one receives
\[
\vect y=\vect c+\vect e \in \F_q^n,
\]
where \(\vect e\in \F_q^n\) is the error vector. Applying the parity-check map yields
\[
H(\vect y)=H(\vect c+\vect e)=H(\vect e),
\]
since \(H(\vect c)=0\) for every \(\vect c\in C\).

Thus decoding may be viewed as follows: from the syndrome \(H(\vect y)\), choose an error estimate \(\widehat{\vect e}\) with the same syndrome, typically of smallest Hamming weight, and define
\[
\widehat{\vect c}=\vect y-\widehat{\vect e}\in C.
\]
Then recover the message by
\[
\widehat{\vect x}=\mathsf{Enc}^{-1}(\widehat{\vect c}).
\]
Hence error correction is the procedure of moving the received word from its ambient coset in \(\F_q^n/C\) back to the code \(C=\ker(H)\), and then inverting the encoding map.

%A linear \([n,k]_q\) code is most naturally described by a short exact sequence
%\[
%0 \longrightarrow \F_q^k \xrightarrow{\;\mathsf{Enc}\;} \F_q^n \xrightarrow{\;\mathsf{Syn}\;} \F_q^{\,n-k} \longrightarrow 0,
%\]
%whose image \(C=\operatorname{im}(\mathsf{Enc})\) is the code and whose kernel satisfies
%\[
%C=\ker(\mathsf{Syn}).
%\]
%If a message \(\vect x\in \F_q^k\) is transmitted, the corresponding codeword is
%\[
%\vect c=\mathsf{Enc}(\vect x)\in C.
%\]
%After the channel introduces an error vector \(\vect e\in \F_q^n\), the received word becomes
%\[
%\vect y=\vect c+\vect e.
%\]
%Applying the syndrome map yields
%\[
%\mathsf{Syn}(\vect y)=\mathsf{Syn}(\vect e),
%\]
%because \(\mathsf{Syn}(\vect c)=0\). Hence the decoder uses the syndrome of \(\vect y\) to determine an error estimate \(\widehat{\vect e}\), then forms
%\[
%\widehat{\vect c}=\vect y-\widehat{\vect e}\in C,
%\]
%and finally recovers
%\[
%\widehat{\vect x}=\mathsf{Enc}^{-1}(\widehat{\vect c}).
%\]
%Thus error correction is the process of returning the received word from its ambient coset in \(\F_q^n/C\) back to the kernel \(C\) of the syndrome map.
\end{note}
\end{flushleft}

\newpage
\subsection{How to Correct Errors}






